Minimally invasive surgery presents a challenging environment for surgical scene understanding because operative video is visually complex, spatially constrained, and clinically contextual. Recent advances in surgical artificial intelligence have enabled progress in tasks such as instrument segmentation, workflow recognition, and surgical action analysis. However, modelling surgical interactions in a clinically meaningful and interpretable way remains difficult. Surgical action triplets, represented as instrument--verb--target relationships, provide a compact description of surgical activity by specifying which instrument performs which action on which anatomical structure. Existing triplet understanding methods largely operate at the frame level, identifying interactions without explicitly grounding them to the responsible instrument instance.

This thesis advances fine-grained surgical scene understanding by developing datasets, benchmarks, and models for grounded instrument--tissue interaction modelling in laparoscopic surgery. First, it reviews the broader landscape of multitask learning and surgical scene understanding in minimally invasive surgical vision, identifying the progression from perceptual and workflow tasks towards fine-grained action understanding, as well as the lack of explicit spatial grounding for surgical interactions.

Second, this thesis introduces CholecInstanceSeg, a large-scale tool instance segmentation dataset for laparoscopic cholecystectomy. CholecInstanceSeg provides instance-level annotations for surgical instruments across routine clinical procedures, enabling individual tools to be localised and distinguished in complex multi-instrument scenes. This dataset establishes the spatial foundation required for grounded surgical action understanding.

Third, this thesis introduces triplet segmentation, a task that extends surgical action triplet recognition from frame-level prediction to instance-level spatial grounding. To support this task, CholecTriplet-Seg is introduced as a benchmark linking instrument instance masks with verb and anatomical target labels. TargetFusionNet is then proposed as a target-aware architecture that integrates weak anatomical priors into an instance segmentation framework for improved grounded triplet prediction. Experiments show that strong instance supervision substantially improves triplet grounding and that target-aware fusion improves full triplet prediction over existing baselines.

Finally, this thesis investigates whether anatomical target prediction should be approached not only as a representation-learning problem, but also as a reasoning problem. Using multimodal vision-language models and structured target-centric reasoning, it explores how visual evidence, anatomical context, ambiguity analysis, and weak priors can support fine-grained surgical interaction understanding.

Collectively, this work moves surgical action understanding beyond frame-level recognition towards grounded, interpretable, and reasoning-aware modelling of instrument--tissue interactions. These contributions provide foundations for future context-aware surgical assistance, workflow analysis, and intelligent computer-assisted intervention systems.
